\begin{statement}[C.1.3; Proved: numerical cross-check]
The independent analytic computation in
\path{code/Cspecial_torsion_numerics.py} gives, to fifty decimals,
\begin{align*}
 N(z_1)&=1.31607401295249246081921890179699905516006859020582\ldots,\\
 N(z_2)&=1.41421356237309504880168872420969807856967187537695\ldots,\\
 \operatorname{Im}\beta^{\rm part}(z_1)
  &=-1.27101206057537346286579673538412682778809507094949\ldots,\\
 \operatorname{Im}\beta^{\rm part}(z_2)
  &=-0.21100425815613411418178314271018934802718590587233\ldots.
\end{align*}
The theta/eta evaluations and exact expressions differ by less than
$10^{-100}$ at working precision $110$ decimal digits.
\end{statement}
\begin{proof}
The program directly evaluates the theta series and the eta product,
then compares the displayed closed forms. The eta product has $160$
factors. Since $|q|\le e^{-\pi\sqrt3}<1/200$, the omitted logarithm is
at most $\sum_{m>160}|q|^m/(1-|q|^m)<10^{-368}$.
The output records the measured floating-point discrepancies. The
exact values are established by C.1.2; this computation is their
independent numerical check, not an additional rounding certificate.
\end{proof}
